\documentclass[aspectratio=169]{beamer}

\usetheme{Madrid}
\setbeamertemplate{navigation symbols}{}
\setbeamertemplate{footline}[frame number]

\usepackage{amsmath, amssymb}
\usepackage{booktabs}
\usepackage{hyperref}

\title{Empirical Designs in Applied Micro}
\subtitle{Shift-Share IV and Unilateral Divorce Laws}
\author{Erika Lynet Salvador}
\date{\today}

\begin{document}

% =====================================================
\begin{frame}
\titlepage
\end{frame}

\begin{frame}{Big Picture}
Three papers, three design lessons:

\begin{itemize}
    \item How to interpret and diagnose shift-share IV.
    \item How staggered policy exposure affects long-run outcomes.
    \item How dynamic treatment effects can overturn prior conclusions.
\end{itemize}

\end{frame}

% =====================================================
\section{Paper 1: Borusyak--Hull--Jaravel (2025, JEP)}

\begin{frame}{Paper 1: Research Question}

What identifies a shift-share (Bartik-style) instrument?

\[
z_i = \sum_{k=1}^K s_{ik} g_k
\]

When is this valid?
Where does quasi-randomness come from?

\end{frame}

\begin{frame}{Paper 1: Core Contribution}

Shift-share IV is not one design.

There are two identification paths:

\begin{enumerate}
    \item Exogenous shifts (\(g_k\))
    \item Exogenous shares (\(s_{ik}\))
\end{enumerate}

Each requires different assumptions and diagnostics.

\end{frame}

\begin{frame}{Paper 1: Identification via Exogenous Shifts}

Interpretation:
The shocks \(g_k\) are quasi-random.

Units differ only in exposure weights.

Key requirement:
\[
\text{Shocks uncorrelated with exposure-weighted unobservables.}
\]

Threat:
Shocks correlated with omitted shift-level fundamentals.

\end{frame}

\begin{frame}{Paper 1: Shift-Based Checklist}

\begin{itemize}
    \item Is the shock plausibly exogenous?
    \item Are there many shocks?
    \item Do shocks predict pre-trends?
    \item Is inference clustered appropriately?
    \item Are incomplete shares handled correctly?
\end{itemize}

\end{frame}

\begin{frame}{Paper 1: Identification via Exogenous Shares}

Interpretation:
Shares behave like valid instruments.

Parallel-trends style condition:

\[
\text{Units with higher } s_{ik}
\text{ would not have differential trends absent treatment.}
\]

Threat:
Shares proxy for long-run structural decline.

\end{frame}

\begin{frame}{Paper 1: Share-Based Checklist}

\begin{itemize}
    \item Do shares satisfy exclusion?
    \item Do pre-trends differ by exposure?
    \item Are results driven by a few dominant shocks?
    \item What do Rotemberg weights look like?
\end{itemize}

\end{frame}

\begin{frame}{Paper 1: Why This Matters}

Shift-share IV requires clarity:

\begin{itemize}
    \item Where does identification come from?
    \item Where does sampling variation live?
    \item What level should inference be done at?
\end{itemize}

\end{frame}

% =====================================================
\section{Paper 2: Gruber (2004, JLE)}

\begin{frame}{Paper 2: Research Question}

Is making divorce easier bad for children in the long run?

Do children exposed to unilateral divorce laws
experience worse adult outcomes?

\end{frame}

\begin{frame}{Paper 2: Methodology}

\begin{itemize}
    \item Exploit staggered state adoption of unilateral divorce.
    \item Compare cohorts exposed vs unexposed in childhood.
    \item Include state and year fixed effects.
    \item Robustness: state trends and migration controls.
\end{itemize}

\end{frame}

\begin{frame}{Paper 2: Outcomes}

Adult outcomes measured years later:

\begin{itemize}
    \item Educational attainment
    \item Family income
    \item Marital stability
    \item Suicide risk
\end{itemize}

\end{frame}

\begin{frame}{Paper 2: Results}

Exposure to unilateral divorce laws in childhood:

\begin{itemize}
    \item ↓ Education
    \item ↓ Income
    \item ↑ Marital instability
    \item ↑ Suicide risk
\end{itemize}

Interpretation:
Family disruption and bargaining changes matter.

\end{frame}

\begin{frame}{Paper 2: Identification Concerns}

\begin{itemize}
    \item Policy endogeneity across states.
    \item Differential trends.
    \item Migration and mismeasurement of exposure.
    \item Mechanism: divorce incidence vs bargaining.
\end{itemize}

\end{frame}

% =====================================================
\section{Paper 3: Wolfers (2006, AER)}

\begin{frame}{Paper 3: Research Question}

Did unilateral divorce laws raise divorce rates?

Re-examines Friedberg (1998).

\end{frame}

\begin{frame}{Paper 3: Methodology}

State-year panel of divorce rates.

Uses event-study dynamics:

\[
\text{DivRate}_{st}
=
\sum_{\ell} \beta_\ell \mathbf{1}\{t-T_s=\ell\}
+ \alpha_s + \lambda_t
\]

Allows effects to evolve over time.

\end{frame}

\begin{frame}{Paper 3: Results}

\begin{itemize}
    \item Large short-run spike in divorce after reform.
    \item Effect fades over time.
    \item Small long-run effect on divorce flows.
\end{itemize}

Conclusion:
Reform accelerates divorces rather than permanently raising them.

\end{frame}

\begin{frame}{Paper 3: Critique of Friedberg}

Single post-reform dummy + trends
can misattribute dynamic effects.

Trend terms may absorb post-spike declines,
creating artificial persistence.

Lesson:
Dynamics matter.

\end{frame}

\begin{frame}{Paper 3: Coasian Interpretation}

If bargaining is efficient:

Legal regime should not greatly affect divorce rates.

Wolfers finds:

\begin{itemize}
    \item Some effect → bargaining frictions exist.
    \item Small long-run effect → frictions limited.
\end{itemize}

\end{frame}

% =====================================================
\section{Synthesis}

\begin{frame}{What These Papers Teach Us}

\begin{itemize}
    \item Identification must be explicit.
    \item Trend specification can drive results.
    \item Dynamic treatment effects matter.
    \item Policy changes operate through outside options and bargaining.
\end{itemize}
\end{frame}

\end{document}
